\def\tituloconcepto{Desacople dinámico}
% Si el capitulo se compila solo (sin main.tex), esta cabecera abre el documento.
% Antes de % Si el capitulo se compila solo (sin main.tex), esta cabecera abre el documento.
% Antes de % Si el capitulo se compila solo (sin main.tex), esta cabecera abre el documento.
% Antes de \input{cabecera} cada capitulo define \tituloconcepto.
\ifdefined\enmain\else
  \documentclass[11pt,a4paper]{article}
  \input{preambulo}
  \begin{document}
  \title{\tituloconcepto}
  \author{\autor}
  \date{\today}
  \maketitle
\fi
 cada capitulo define \tituloconcepto.
\ifdefined\enmain\else
  \documentclass[11pt,a4paper]{article}
  % Preambulo compartido por todos los documentos. Compilar con XeLaTeX/LuaLaTeX
% (o Tectonic). Si se usa Overleaf, elegir el compilador XeLaTeX.
\usepackage[spanish,es-noshorthands]{babel}
\usepackage{amsmath,amssymb}
\usepackage{booktabs,array}
\usepackage[a4paper,margin=2.5cm]{geometry}
\usepackage{xcolor}
\usepackage{listings}
\usepackage[most]{tcolorbox}
\usepackage{quantikz}
\usepackage[hidelinks]{hyperref}

% ---- Autor (editar a gusto) -------------------------------------------------
\newcommand{\autor}{Genesis Silesky}

% ---- Notacion de Dirac y atajos ---------------------------------------------
% quantikz ya define \ket y \bra; DeclareRobustCommand los redefine sin error.
\DeclareRobustCommand{\ket}[1]{\left|#1\right\rangle}
\DeclareRobustCommand{\bra}[1]{\left\langle#1\right|}
\DeclareRobustCommand{\braket}[2]{\left\langle#1\middle|#2\right\rangle}
\newcommand{\C}{\mathbb{C}}
\newcommand{\R}{\mathbb{R}}
\DeclareMathOperator{\Tr}{Tr}

% ---- Cajas de apoyo --------------------------------------------------------
% "concepto": definiciones previas que hacen falta para entender el capitulo.
\newtcolorbox{concepto}[1]{breakable,colback=blue!4,colframe=blue!55!black,
  fonttitle=\bfseries,title={Concepto previo: #1}}
% "nota": aclaraciones y advertencias.
\newtcolorbox{nota}[1]{breakable,colback=orange!6,colframe=orange!70!black,
  fonttitle=\bfseries,title={#1}}

% ---- Estilos de listados ----------------------------------------------------
\lstdefinestyle{py}{language=Python,basicstyle=\ttfamily\footnotesize,
  breaklines=true,postbreak=\mbox{\textcolor{gray}{$\hookrightarrow$}\space},
  numbers=left,numberstyle=\tiny\color{gray},frame=single,
  commentstyle=\color{green!40!black}\itshape,
  keywordstyle=\color{blue!70!black}\bfseries,
  stringstyle=\color{red!60!black},showstringspaces=false,tabsize=4,
  columns=fullflexible,keepspaces=true}
\lstdefinestyle{salida}{basicstyle=\ttfamily\footnotesize,breaklines=true,
  frame=leftline,backgroundcolor=\color{gray!8},columns=fullflexible,
  keepspaces=true}

  \begin{document}
  \title{\tituloconcepto}
  \author{\autor}
  \date{\today}
  \maketitle
\fi
 cada capitulo define \tituloconcepto.
\ifdefined\enmain\else
  \documentclass[11pt,a4paper]{article}
  % Preambulo compartido por todos los documentos. Compilar con XeLaTeX/LuaLaTeX
% (o Tectonic). Si se usa Overleaf, elegir el compilador XeLaTeX.
\usepackage[spanish,es-noshorthands]{babel}
\usepackage{amsmath,amssymb}
\usepackage{booktabs,array}
\usepackage[a4paper,margin=2.5cm]{geometry}
\usepackage{xcolor}
\usepackage{listings}
\usepackage[most]{tcolorbox}
\usepackage{quantikz}
\usepackage[hidelinks]{hyperref}

% ---- Autor (editar a gusto) -------------------------------------------------
\newcommand{\autor}{Genesis Silesky}

% ---- Notacion de Dirac y atajos ---------------------------------------------
% quantikz ya define \ket y \bra; DeclareRobustCommand los redefine sin error.
\DeclareRobustCommand{\ket}[1]{\left|#1\right\rangle}
\DeclareRobustCommand{\bra}[1]{\left\langle#1\right|}
\DeclareRobustCommand{\braket}[2]{\left\langle#1\middle|#2\right\rangle}
\newcommand{\C}{\mathbb{C}}
\newcommand{\R}{\mathbb{R}}
\DeclareMathOperator{\Tr}{Tr}

% ---- Cajas de apoyo --------------------------------------------------------
% "concepto": definiciones previas que hacen falta para entender el capitulo.
\newtcolorbox{concepto}[1]{breakable,colback=blue!4,colframe=blue!55!black,
  fonttitle=\bfseries,title={Concepto previo: #1}}
% "nota": aclaraciones y advertencias.
\newtcolorbox{nota}[1]{breakable,colback=orange!6,colframe=orange!70!black,
  fonttitle=\bfseries,title={#1}}

% ---- Estilos de listados ----------------------------------------------------
\lstdefinestyle{py}{language=Python,basicstyle=\ttfamily\footnotesize,
  breaklines=true,postbreak=\mbox{\textcolor{gray}{$\hookrightarrow$}\space},
  numbers=left,numberstyle=\tiny\color{gray},frame=single,
  commentstyle=\color{green!40!black}\itshape,
  keywordstyle=\color{blue!70!black}\bfseries,
  stringstyle=\color{red!60!black},showstringspaces=false,tabsize=4,
  columns=fullflexible,keepspaces=true}
\lstdefinestyle{salida}{basicstyle=\ttfamily\footnotesize,breaklines=true,
  frame=leftline,backgroundcolor=\color{gray!8},columns=fullflexible,
  keepspaces=true}

  \begin{document}
  \title{\tituloconcepto}
  \author{\autor}
  \date{\today}
  \maketitle
\fi


\section{Desacople dinámico: apagar el acoplamiento con pulsos $\pi$}

\subsection{Objetivo}
Ver cómo el acoplamiento entre dos espines hace que la fase de uno dependa del
estado del otro, y cómo un pulso $\pi$ a mitad de la evolución lo cancela
(\emph{desacople}).

\subsection{Conceptos previos}
\begin{concepto}{acoplamiento $\sigma_z\otimes\sigma_z$ y spin echo}
Ver capítulo~12 (el acoplamiento $H=2\pi J\,\sigma_z\otimes\sigma_z/4$ y cómo el
eje de giro de un qubit depende del otro) y capítulo~21 (un pulso $\pi$ revierte
las fases acumuladas).
\end{concepto}

\subsection{Fundamento teórico}
El desacople dinámico es una técnica de control que usa pulsos $\pi$ para
suprimir interacciones no deseadas, y la idea es la misma del spin echo.

\paragraph{Modelo.} Se usan dos espines: $H$ (un protón) y $P$ (un núcleo de
fósforo), con acoplamiento $J=696$~Hz (el valor del capítulo~12). Durante un
tiempo libre $t$,
\begin{equation}
  U(t)=\exp\!\big(-i\,2\pi J\,\sigma_z\otimes\sigma_z\,t/4\big).
\end{equation}
Si $H$ está en $\ket0$ (arriba), $P$ gira alrededor de $z$ con una frecuencia, y si
está en $\ket1$ (abajo) con la contraria. La fase que acumula $P$ es
$\pm\pi Jt$.

\paragraph{Superposición.} Si $H$ está en una superposición, $P$ gira en los
dos sentidos a la vez, y los dos espines se entrelazan: el vector de Bloch de
$P$ se acorta ($|\cos\pi Jt|$ para la superposición $50/50$) y su fase se
pierde.

\paragraph{Desacople con un pulso $\pi$.} Se deja evolucionar $t/2$, se aplica un
pulso $\pi$ y se deja evolucionar otros $t/2$:
\begin{itemize}
  \item \textbf{$\pi$ sobre $H$:} invierte el estado de $H$ y, con él, el signo
    del giro de $P$ en la segunda mitad; las fases se cancelan.
  \item \textbf{$\pi$ sobre $P$:} invierte la fase de $P$ en la mitad; lo que
    acumula después deshace lo anterior (como en el spin echo).
  \item \textbf{$\pi$ sobre ambos:} ambos efectos se suman en vez de cancelarse.
    El producto $\sigma_z\otimes\sigma_z$ no cambia al invertir los dos espines,
    así que el acoplamiento sigue activo.
\end{itemize}
Se considera que hubo desacople si $P$ termina igual que lo haría sin acoplamiento
($J=0$) con la misma secuencia de pulsos.

\begin{nota}{Fuentes y alcance}
Este capítulo se basa en la guía de la práctica de desacople dinámico y en el
modelo de acoplamiento del capítulo~12. El archivo de teoría de esta práctica es
un escaneo sin texto legible, así que las conclusiones salen \textbf{de calcular}
el modelo, no de citar el documento. No se cubre el último ejercicio de la
guía (diseñar un conjunto de pulsos para 5 espines con una matriz de
desacople).
\end{nota}

\subsection{Código (completamente comentado)}
\lstinputlisting[style=py]{../../fuentes/22_desacople_dinamico.py}

\subsection{Explicación celda por celda}
\begin{description}
  \item[Celda 2] Herramientas: evolución libre, pulso $\pi$ sobre $H$, $P$ o
    ambos, vector de Bloch de $P$ (traza parcial) y su fase.
  \item[Celda 3] Fase de $P$ con $H$ en $\ket0$ y en $\ket1$ para $t=0.5$, $1$ y
    $2$~ms.
  \item[Celda 4] Longitud transversal de $P$ con $H$ en superposición.
  \item[Celda 5] Aplica un pulso $\pi$ a mitad de la evolución y compara la fase de $P$
    con y sin acoplamiento.
\end{description}

\subsection{Salida obtenida}
\lstinputlisting[style=salida]{../../fuentes/salidas/22_desacople_dinamico.txt}
Es determinista.

\subsection{Interpretación}
La fase de $P$ vale $\pm\pi Jt$: $62.6^\circ$ a $0.5$~ms y $125.3^\circ$ a $1$~ms,
con signo opuesto según el estado de $H$. Con $H$ en superposición $50/50$, la
longitud transversal de $P$ baja a $0.460$ ($0.5$~ms) y a $0.578$ ($1$~ms),
que coincide con $|\cos\pi Jt|$. Con un pulso $\pi$ sobre $H$ o sobre $P$ a
mitad de la evolución, la fase de $P$ es $0^\circ$, igual que sin acoplamiento:
\textbf{hay desacople}. Con $\pi$ sobre ambos, la fase queda en $\mp125.3^\circ$
(a $1$~ms): \textbf{no hay desacople}.

\subsection{Ejercicios}
\begin{enumerate}
  \item Cambia la posición del pulso $\pi$ (por ejemplo, a $t/4$) y observa si
    hay desacople.
  \item Con $H$ en superposición $50/50$, aplica un pulso $\pi$ sobre $H$ a mitad
    y mide la longitud transversal de $P$.
  \item Prueba con $J=300$~Hz y con $t=2$~ms: ¿cambia la conclusión?
\end{enumerate}

\subsection{Errores comunes}
\begin{itemize}
  \item Creer que invertir los dos espines ayuda: como $\sigma_z\otimes\sigma_z$
    no cambia, el acoplamiento sigue.
  \item Olvidar comparar con el caso $J=0$: el pulso $\pi$ también cambia la
    fase de $P$ por sí mismo.
\end{itemize}

\ifdefined\enmain\else
  \end{document}
\fi

